
\subsubsection{Experimental Framework}

We structure our investigation as four sub-hypotheses tested sequentially: (h-e1) transfer-stable categories exist, (h-m1) optimal low-level thresholds transfer across stages, (h-m2) objective-dependence predicts categorical separation, and (h-m3) embedding-stage mismatch creates quality-speed trade-offs. Each hypothesis uses controlled three-way comparisons (baseline, transferred, stage-tuned) to isolate curation effects.

\textbf{h-e1 (Existence).} Low-level quality filters (deduplication, perplexity) applied with pre-training-derived thresholds to fine-tuning data will perform within 1\% of stage-tuned thresholds on MMLU and HellaSwag. This establishes that not all curation requires stage-specific optimization.

\textbf{h-m1 (Mechanism - Universal Hygiene).} Optimal thresholds for low-level operations will not vary significantly (>10\% performance delta when mismatched) across pre-training and fine-tuning, because these operations address data quality properties independent of training objectives.

\textbf{h-m2 (Mechanism - Categorical Separation).} Curation techniques categorized by objective-dependence will show distinct transfer profiles: objective-independent techniques $\leq 1$\% delta, objective-dependent >5\% degradation. Requires non-overlapping confidence intervals to confirm categorical separation.

\textbf{h-m3 (Mechanism - Quality-Speed Trade-Off).} Embedding-based subset selection using early-stage embeddings (MiniLM) versus late-stage embeddings (Instructor) will exhibit measurable trade-offs: stage-mismatch penalty >2\%, compute cost ratio >$3\times$.

\subsubsection{Datasets and Models}

\textbf{Pre-training data:} C4 52,002-sample subset. Standard web-scraped corpus with documented filtering (dedup threshold 0.8, perplexity ~1000).

\textbf{Fine-tuning data:} Dolly-15k (15,000 samples, split 13,509 train / 1,502 validation) and Alpaca-52k (52,002 samples). Open-domain instruction-response pairs.

\textbf{Model:} Llama-2-7B. Mid-size causal language model balancing experimental feasibility with meaningful performance measurement.

\textbf{Evaluation:} MMLU (5-shot) and HellaSwag (10-shot). Standard benchmarks sensitive to 1-2\% performance differences.

\subsubsection{Curation Techniques}

\textbf{Objective-Independent (Low-Level):}

\textit{Deduplication.} MinHash LSH with Jaccard similarity threshold. Removes near-duplicate samples. Operates on surface n-gram statistics independent of task objective. Threshold range: 0.7-0.9 (C4 uses 0.8).

\textit{Perplexity filtering.} Removes samples with high language model perplexity (gibberish, foreign language, corruption). \textbf{PoC limitation:} Used length-based proxy (not GPT-2 or KenLM); filtered 0 samples, limiting perplexity validation to infrastructure only. Threshold range: 500-1500 (C4 uses ~1000).

\textbf{Objective-Dependent (High-Level):}

\textit{Instruction quality filters.} Task-specific heuristics for instruction clarity, response quality, prompt diversity. Directly depend on instruction-following objective. For Dolly, we use prompt diversity (lexical similarity < 0.5) as proxy. True multi-source domain mixing requires datasets with source metadata (unavailable in Dolly).

\subsubsection{Controlled Variables}

Model architecture (Llama-2-7B), hyperparameters (lr=2e-5, batch=8, epochs=3), evaluation protocol (fixed 5-shot MMLU, 10-shot HellaSwag), and pre-training checkpoint held constant across conditions. Independent variable: curation threshold source (C4 pre-training vs. Dolly/Alpaca fine-tuning-optimized).

\subsubsection{Statistical Validation}

Bootstrap 95\% confidence intervals with 10,000 resamples. Welch's t-test comparing objective-independent vs. objective-dependent transfer deltas. Cohen's d for effect size (threshold d > 0.8 indicates large effect). Significance level p < 0.05 for primary claims, p < 0.10 acceptable for secondary hypotheses.

\subsubsection{Proof-of-Concept Execution}

Given resource constraints, we execute at PoC tier with mock evaluation: curation pipelines run on full datasets, but model training is simulated using predicted scores from published benchmarks (Llama-2 technical report, DataComp results). This validates mechanism direction and infrastructure while deferring absolute precision to production.

\textbf{PoC simplifications:} Mock MMLU/HellaSwag scores (not actual lm-eval), length-based perplexity proxy (not KenLM), exact-match deduplication (not fuzzy LSH), single seed (not multi-seed statistical robustness).

\textbf{PoC validates:} Transfer delta patterns, categorical separation direction, quality-speed trade-off ratios. Absolute performance values subject to production verification.
